\begin{afterword}
\summaryhead

The post starts from Carl Sagan's parable of someone who claims that an invisible dragon lives in their garage and, as each test is proposed, explains why it will fail: the dragon is inaudible, does not breathe, is permeable to flour. Sagan used the story to show how a poor hypothesis dodges falsification. The author draws another point: to know in advance which tests to excuse, the claimant must expect, somewhere in their mind, that there is no dragon.

The explanation offered is Daniel Dennett's ``belief in belief'': it is often easier to believe that you ought to believe something than to believe it. The author adds that a person who only believes in belief cannot admit this to themselves, and that belief in belief is often wordless, a flinch from admitting error or a wish to protect one's self-image. So the concept of belief must include unspoken anticipations. Excuses made in advance are the sign that belief and belief in belief have come apart.

\responsehead

The post's central observation is its own, and it is good. Sagan concluded from the excuses that ``something funny is going on inside my head,'' and asked what had convinced the claimant. The post sees that the excuses are aimed exactly at the tests that would fail, so the claimant's expectations must be accurate. That is a real inference, and it gives the post its one observable sign: excuses made in advance.

The rest of the post sets itself against earlier thinkers who, as far as I could check, had reached similar ideas. The unnamed philosophers who ask whether the claimant ``really'' believes are dismissed with a joke about disk space. But their question is the standard puzzle of self-deception: how can a person hold contradictory beliefs at once? The answers proposed before 2007 include Robert Audi's, that the self-deceived sincerely avow what they do not believe, which is close to the post's own answer. The post also presents its wider concept of belief, including wordless anticipation, as a step beyond Dennett. Dennett had drawn a similar line himself in 1978, between beliefs that animals share and verbal ``opinions,'' and had located self-deception where the two diverge. I had this from a secondary account and could not read Dennett's chapter.

The psychology itself is asserted, not shown. ``Nobody will admit to themselves'' that they only believe in belief; wordless states ``account for as much of ourselves as words''; one account of the claimant is ``not psychologically realistic'' and another is. No study, case or observation beyond the parable, the Santa Claus memory and the reader's own garage is offered for any of these. They may be true, but the post gives the reader no way to check them except introspection.

In short: a sharp inference from Sagan's dragon, joined to a psychology of belief that is asserted without evidence and set against unnamed philosophers and against Dennett, some of whom, as far as I could check, had drawn similar distinctions earlier.
\end{afterword}
